\chapter{Conclusion} \label{ch:conclusion}

The research was initiated with the growing demand for a computational tool that is able to simulate large chemical reaction networks and to predict pollutant emission. We applied the Julia programming language to solve the one- and the two-step mechanism of methane combustion, the 5-node-5-edge networks and larger line-shaped networks. 

The results of performance measurement prove that Julia with its ecosystem is competent in simulating large network dynamics within reasonable time as long as appropriate numerical solvers are chosen. However, unlike C/C++ and Python, Julia has a rather shorter history. The Julia core code and its packages are still under heavy development thanks to the common effort from the open-source Julia community. It can be foreseen that Julia will be shortly utilized to speed up CRN simulations in industry and to estimate emissions more accurately for industrial furnaces.